\documentclass[11pt,a4paper]{article}
\usepackage[T1]{fontenc}
\usepackage[utf8]{inputenc}
\usepackage{lmodern,microtype}
\usepackage[margin=.78in]{geometry}
\usepackage{amsmath,amssymb,mathtools,booktabs,tabularx,array,enumitem}
\usepackage[hidelinks]{hyperref}
\usepackage{fancyhdr}
\setlength{\headheight}{14pt}
\pagestyle{fancy}\fancyhf{}
\lhead{\small Sync HotStuff RR, Part III}\rhead{\small Staking / v1}
\cfoot{\thepage}
\setlength{\parindent}{0pt}\setlength{\parskip}{5pt}
\setlength{\emergencystretch}{2em}
\setlist{nosep,leftmargin=*}
\newcolumntype{Y}{>{\raggedright\arraybackslash}X}
\newcommand{\Hash}{\mathsf H}
\newcommand{\tagg}[1]{\mathsf{#1}}
\newcommand{\cfg}{\mathcal C}
\newcommand{\sid}{s}
\hypersetup{pdftitle={Sync HotStuff RR, Part III: Staking and Authority, v1}}
\begin{document}
\begin{center}
{\Large\bfseries Part III: Staking and Authority}\\[3pt]
{\small Version 1 \quad 19 October 2026}
\end{center}

This module turns locked stake into a fixed regional committee and a leader
for each view. L1 updates stake in long epochs. Regional sessions may rotate
much faster: several sessions can use one snapshot. The next session's
committee and base leader schedule are known when its predecessor starts.
Regular handoff still needs no new L1 transaction.

Withdrawal uses one additional control field: a minimum snapshot number.
An effective multisignature Force can raise this minimum while closing
the old branch. This gives a concrete end to all old duties, including
speculative sessions. Final withdrawals therefore depend on a batch Force;
the cost and timing are stated in Section~5.

\section{Objects and the division of work}
\label{sec:objects}
Part II supplies a regional BFT engine, regular handoff and Force. A
\emph{session} $s=(\beta,r)$ is one region's production turn. Here $\beta$
identifies the branch and $r$ the turn, not the physical region. Its views
are $(\beta,r,g,t)$: recovery increases $g$, and an ordinary leader change
increases $t$. The initial view is $(\beta,r,0,0)$.

Write $A_s$ for the session's initial execution anchor and $T_s$ for its
terminal DATA. Its fixed quota $\mathrm{TB}_s$ gives the terminal height
$h(A_s)+\mathrm{TB}_s$. A \emph{FINAL receipt} certifies regional commitment;
a proposal QC alone does not. Regular handoff continues from $T_s$ after
its complete receipt. Espresso-finalized history and the finalized L1
controls determine the accepted global prefix $G$. An effective Force
closes the whole current branch and restarts from the retained tip of $G$.
It cannot remove anything already in $G$.

\begin{tabularx}{\linewidth}{@{}p{.24\linewidth}Y@{}}\toprule
Object & Fields and purpose\\\midrule
Stake position $o$ & Unique ID, region, operator key, withdrawal owner,
remaining escrow $B_o$, queue state, exit epoch, and processed offenses.\\
Snapshot $\cfg_e$ & Epoch $e$, fixed cutoff, future activation height $a_e$,
members and integer weights, and a binding configuration ID.\\
Branch policy & Multisig-fixed candidate-order rule, ON/OFF map $S_\beta$,
and minimum snapshot number $m_\beta$.\\
Session plan & Target pair, selection reference $E_s$, configuration,
region, and base leader rule. It permits preparation, not signing.\\
Session authority & The plan, exact execution anchor $A_s$, handoff
evidence, quota and regional BFT parameters.\\\bottomrule
\end{tabularx}

One operator key belongs to at most one live position; IDs are never
reused. This first version uses one token and whole-position joins and
exits. Key or region changes require a new registration after the old
liability is settled. Historical records remain available.

L1 keeps money, queues, snapshots, controls and penalties. Nodes derive
session authority from authenticated history. A sound replay verifier
can expose the same history to the withdrawal contract; Section~5 states
exactly what it must verify. L1 need not track every regular session.

\newpage
\section{Three public selection interfaces}
\label{sec:selection}
\paragraph{Snapshot lookup.}
At finalized L1 reference height $\ell$, let
\[
 e(\ell)=\max\{e:a_e\le\ell\}.
\]
An existence proof for an old snapshot is insufficient: the reader checks
the complete applicable registry history. Missing records or member data
cause a wait. They do not justify using an older snapshot.

For a regular successor $s=(\beta,r+1)$ of $p=(\beta,r)$, set
\begin{equation}\label{eq:reference}
 E_s=A_p,\qquad h_s^{\rm sel}=h(E_s),\qquad
 e_s=\max\{m_\beta,e(\lambda(E_s))\}.
\end{equation}
The selected weights are those of $\cfg_{e_s}$. Normally the minimum is
inactive, so this is exactly the latest snapshot at the predecessor
anchor's L1 reference. After a withdrawal Force, the minimum prevents an
old anchor from reauthorizing released stake. It is a branch-fixed rule,
not a query for the latest local L1 head.

\paragraph{Region order.}
The multisig fixes a deterministic rule for
$\tagg{randList}(\beta,r)$, a permutation of the registered regions used
for $p$'s regular successor. Choose the first region $i$ for which
$S_\beta(i)=\tagg{ON}$ and $W_i=\sum_{o\in V_i(\cfg_{e_s})}w_o>0$.
If none qualifies, the boundary is pending. Missing evidence is not zero
stake, and silence does not authorize skipping a selected region. The
policy specifies whether the current region is last. It also supplies a
separate ordered boot list for the first session of a new branch.

\paragraph{Weighted choice.}
$\tagg{WeightedRand}(\mathcal W,z)$ is a pure function. Sort positive
integer weights by position ID, derive an integer $x\in[0,W)$ from $z$
using specified hash expansion and rejection sampling, and select the
interval containing $x$. For weights $(60,30,10)$, $x=73$ selects the
second member. An empty set returns $\bot$.

The session's base leader function is
\begin{equation}\label{eq:leader}
 L_s^0(g,t)=\tagg{WeightedRand}\bigl(\mathcal W_s,
 \Hash(\texttt{LEADER},\mathit{chainId},h_s^{\rm sel},\beta,r_s,g,t)\bigr).
\end{equation}
Here $r_s$ is the \emph{target} turn, and $\mathcal W_s$ contains only its
selected region. Fixed-length or length-prefixed encoding, the hash
algorithm and ID order are protocol constants. Given the inputs,
every node obtains one leader; different views may select the same one.

The hash uses a public DATA height, not transaction contents. With a
uniform sampling input the selection probability is $w_o/W_i$. Public
height hashing alone does not establish unpredictability, resistance to
stake/identity selection, or schedule fairness. We make no such claim.
Section~4 explains how finalized penalties restrict future proposer
eligibility without changing the session's voting weights.

\paragraph{Example.}
For $(0,0)\to(0,1)\to(0,2)$, $E_{(0,2)}=A_{(0,1)}=T_{(0,0)}$.
If that block has DATA height 100 and L1 reference 2500, the base seed is
100 and the snapshot is $\cfg_{\max(m_0,e(2500))}$. The list is
$\tagg{randList}(0,1)$, while the hash names target $(0,2,g,t)$.
If $(0,1)$ ends at DATA 120, $(0,2)$ executes after 120, not after 100.

\newpage
\section{Long epochs and stake queues}
\label{sec:queues}
Choose a long L1-height epoch length $M$ and an advance interval $P$,
with $0<P<M$. Snapshot $e$ activates at $a_e=a_0+eM$ and has cutoff
$c_e=a_e-P$. Genesis supplies $\cfg_0$. Each later cutoff reads the
complete L1 state strictly before block $c_e$, applies the queue rules
below, and freezes $\cfg_e$. Its publication must precede $a_e$.
Late materialization cannot move the cutoff or reorder requests. If a
required snapshot has not been published and verified, its consumers wait.

This separates three events: a request is accepted, a future snapshot is
fixed, and a session adopts it. L1 does not need to know which session is
running at any of these events. Updating $\cfg_e$ never edits a running
session or an already fixed plan.

\paragraph{Join.}
Anyone may deposit, register a region and operator key, and prove key
possession and withdrawal-owner authorization. The position enters its
region's inbound FIFO in L1 order. It has no duties until included in a
snapshot and selected for a session. A fixed entry delay must have elapsed
at the cutoff. Once included in a frozen snapshot, its escrow cannot be
refunded through the pending-request cancellation path.

\paragraph{Exit.}
The owner may enqueue one irrevocable whole-position exit. Until selected
at a cutoff after the fixed exit delay, the position stays eligible.
Selection records an exit epoch $e_x$: the position is absent from every
$\cfg_e$ with $e\ge e_x$. It must still serve old assignments. An exit
request, omission from a snapshot, or elapsed epoch does not itself
authorize withdrawal. The final condition is in Section~5.

\paragraph{Bounded changes.}
Each region has positive inbound and outbound stake budgets and request
count limits per epoch. Select each FIFO's longest mature prefix fitting
its budgets. Require every admitted whole position to fit a fresh budget,
so an oversized head cannot block the queue forever. Remove canceled and
completed entries. Count inflow and outflow separately, even if their net
sum is zero. No regional stake floor may veto an exit indefinitely.

At the cutoff, included positions receive their remaining positive
balances as weights. Positions below the minimum bond are omitted and
retired from future snapshots; associated queue requests are consumed.
Penalty losses and these mandatory removals are not charged to voluntary
exit budgets. A later penalty reduces escrow, not a frozen denominator.
New identities cannot reuse an exited position's funds before release.

\paragraph{Membership and responsibility.}
Use separate records for escrow and duties. A position may be in a
running committee and a prepared successor at the same time. All selected
voters, not just the leader, must remain available. The node retains or
durably transfers Part II's required data and outstanding sends; stopping
a program does not erase that obligation.

These queues borrow Ethereum's separation of pending deposits, exit
epochs and withdrawability, not its waiting constants. They do not prove
the regional fault bound. For each fixed session, Part II still assumes
Byzantine weight at most $F_i<W_i/2$, with certificate threshold $>F_i$,
including key compromises during that authority's accepted lifetime.
Queue progress also requires continued L1 processing and positive budgets.

\newpage
\section{Slashing and recovery leaders}
\label{sec:slash}
All signed evidence binds the chain, protocol, historical position and
configuration, region, session, and its initial execution anchor. View
messages also bind their view and position. Signatures are checked against
that historical authority, never the current balance.

\begin{tabularx}{\linewidth}{@{}p{.25\linewidth}Y@{}}\toprule
Offense & Two signatures of the same authorized key establish\\\midrule
Conflicting proposals & Different admissible DATA/DRAIN proposals at the
same $(\beta,r,g,t,j)$, with Part II's admission and quota evidence.\\
Conflicting fresh votes & Different block choices at the same
$(\beta,r,g,t,j)$; DATA and DRAIN share this one-vote restriction.\\
Conflicting NEWVIEW & Different openings for one $(\beta,r,g)$, even if
the openings are ancestor and descendant.\\
Conflicting ANCHOR votes & Different openings voted for in one
$(\beta,r,g)$.\\
Conflicting FINAL & Different DATA hashes at the same height within one
session; FINAL itself has no generation or term.\\\bottomrule
\end{tabularx}

Timeouts, BLAME, OFF status and forwarding evidence are not these offenses.
Different QCs for the same opening are not NEWVIEW equivocation. Reusing a
height on a different branch is not by itself a conflict.

\paragraph{Evidence and payment.}
Anyone may submit a complete proof to L1. The contract verifies the
historical authorization, signatures and offense predicate. A valid proof
atomically records the offense, debits a published penalty fraction of
the historical bonded weight capped by remaining escrow, and permanently
marks the key ineligible for future proposer records. Invalid or incomplete
proofs do not freeze funds. Deduplicate by position, offense type and
conflict slot, not by the serialization of an evidence pair. Total debits
cannot exceed escrow. Timely valid offenses remain payable during exit;
the first version needs no separate discretionary dispute tribunal.

\paragraph{Recovery record.}
To make exclusion concrete without changing voting weights, retain L1 on
the recovery path. For $g>0$, anyone may submit the complete quit proof
for $(\beta,r,g-1)$ and its authority dependencies. The contract creates
at most one immutable record
\[
 \tagg{REC}(\beta,r,g)=(\mathit{authorityId}_s,g,J_g),
\]
where $J_g$ is exactly the banned-key set in the contract immediately
before that transaction. The proof must justify entry to $g$; a bare
counter cannot reserve arbitrarily many future records. The next
generation waits for this record's finality, as Part II already waits for
recovery authority. $J_g$ changes proposer eligibility only. The leader
is computed by (\ref{eq:leader}) on the remaining members with their
original weights. No eligible proposer means wait for external recovery.

Initial generation $g=0$ uses the proposer exclusions recorded at the
selected snapshot's cutoff. A later slash does not rewrite an already
fixed initial schedule. New snapshots and later recovery records apply
the ban. Old QCs and FINAL receipts retain their original voting weights.

This implements eventual exclusion if valid proofs reach L1 and later
recovery records are created. It does not prove fairness of the public
hash schedule or replace Part II's uninterrupted-service conditions.
Safety and conditional progress keep their separate assumptions.

\newpage
\section{A concrete withdrawal boundary}
\label{sec:withdraw}
An exit cannot be cleared by examining $G$ alone: a region may have
prepared or run later sessions whose DATA has not reached $G$. This
version clears exits in batches through an existing multisig Force.

\paragraph{The minimum snapshot.}
Extend the L1 branch-update record with a candidate-order policy and a
requested minimum $m$. Along the valid control log, minima never decrease;
the requested snapshot must already be activated at the update's L1
publication height. On a matching update's Espresso activation, Part II
closes the old branch and installs the new policy, status map and minimum.
A stale update changes none of them. Genesis sets $m_0=0$.

Every new authority, including a Force bootstrap, obeys
$e_s\ge m_\beta$ through (\ref{eq:reference}). Every later branch inherits
or raises the minimum. Hence a position absent from all snapshots starting
at $e_x$ acquires no new duty after a valid Force installs $m\ge e_x$.
The same Force closes all older running and prepared sessions, including
those ahead of $G$. It does not require enumerating them on L1.

\paragraph{Force bootstrap.}
The first session of a new branch uses the retained tip of $G$ for both
$E_s$ and $A_s$, the new policy's boot list, and
$\cfg_{\max(m_\beta,e(\lambda(E_s)))}$. The target $r$ is Part II's target,
including its zero-output reuse rule. Subsequent regular sessions use
their predecessor's initial anchor again. No old plan is moved to the
new branch, even if a turn number is reused.

\paragraph{What the withdrawal contract verifies.}
For an exiting position, supply its exit record and proof of an effective
Force whose installed minimum is at least $e_x$. From an authenticated
replay checkpoint, verify the complete applicable L1 controls and
Espresso history through that Force's activation, including the fact that
its target was current. This is a proof of the public Part II transition,
not mere Espresso inclusion of a claimed result. Full replay or a sound
proof of that execution implements the same check. Existing L1 and
Espresso trust anchors are assumed; exiting regional keys do not certify
their own retirement. The accepted checkpoint records $(q,G,\beta,m)$.

\paragraph{Time for evidence.}
Start the holding period when L1 accepts that verified boundary, at time
$t_c$. Part II bounds closure visibility by $T_E$ and delivery of an
available complete regional result by $T_{\rm pub}+T_E$. Let $T_{\rm ev}$
additionally bound checking, collecting and submitting relevant evidence
after those deliveries, including evidence about discarded sessions.
Extend the L1 service assumption to valid slash transactions with bound
$T_{\rm L1}^{\rm stake}$. A conservative choice is
\[
 T_{\rm hold}\ \ge\ 2T_E+T_{\rm pub}+T_{\rm ev}
                 +T_{\rm L1}^{\rm stake}+\varepsilon_{\rm clock}.
\]
The deployed L1 timelock must not expire earlier than this real-time
interval; its clock conversion and margin are explicit deployment
parameters. The existing bounds alone do not supply $T_{\rm ev}$.

Release only remaining escrow strictly after the deadline, after all
valid proofs accepted by that deadline have been applied in L1 order.
No fixed window catches evidence hidden forever or unbounded offline
investigation. The release guarantee is conditional on the stated
evidence service, not a claim that every offense will eventually be paid.

\newpage
\section{Integration, guarantees and deliberate costs}
\label{sec:integration}
\paragraph{Regular startup.}
When $p=(\beta,r)$ starts, nodes derive the plan for $(\beta,r+1)$ from
$A_p$, the fixed branch policy and snapshot rule. When $T_p$'s complete
receipt arrives, they bind the actual anchor, verify the handoff and
prepare its execution state. Part II's Ready and Enter then start the
new program. A prediction, a plan or a QC alone never permits entry.

Authorities are now keyed by $(\beta,r)$ and explicit configuration ID,
not by $r$ alone. This replaces Part II's old requirement to reuse one L1
authority whenever Force reuses $r$. The handoff checker, FINAL checker
and global replay all derive the same new context. For the first DATA,
the L1 reference floor includes the incoming anchor's reference, the
selected snapshot's publication height and the branch-control publication
height; it must also reach the snapshot's activation height. Quota and
voting parameters remain fixed within the session.

\paragraph{Why the rules are deterministic.}
The authenticated reference and branch minimum identify one snapshot.
The fixed list and status map identify one eligible region or no region.
Canonical weighted selection and the complete view identify one base
leader. A finalized recovery record fixes the only allowed exclusion set
for its generation. None of these choices uses receipt arrival order or
a reader's latest local head. Missing dependencies cause a wait.

\paragraph{Why voluntary exit does not erase existing liability.}
Before the verified boundary, the contract keeps the position locked.
At the boundary, every older branch is closed and every future authority
uses an epoch at least $e_x$, which excludes this position. Thus no old
program or future Force can extend its accepted authority. Escrow remains
available throughout the subsequent evidence interval. Later penalties
do not retroactively change certificates, and earlier penalties are not
counted again as money still available.

\paragraph{External finality.}
Accepted history is anchored in finalized Espresso and L1 inputs, not in
a later vote by the exited regional committee. This module adds no social
checkpoint mechanism for regional history. It retains the external trust
anchors and sound replay verification: raw Espresso inclusion alone is
not acceptance into $G$. The minimum and branch closure prevent an old
key from reopening authority after withdrawal. Data needed for historical
validation must still be retained or transferred under Part II.

\paragraph{Costs that this version accepts.}
\begin{itemize}[itemsep=2pt]
\item Public schedules are predictable; they can expose future leaders
to targeting and do not prove cryptographic election fairness.
\item A fast regular rotation can reuse a snapshot. A slow session can
delay adoption. The withdrawal Force provides a finite boundary only
when the multisig and required external services actually act.
\item Batch withdrawal can discard the speculative suffix beyond $G$.
It is not a hidden action at every staking epoch or regular handoff.
\item Recovery retains an L1 wait. A session's stake can be slashed below
its frozen weight; its BFT fault assumption is still about that weight.
\end{itemize}

\paragraph{References.}
Part II v4.5, especially fixed snapshots, FINAL, Force-close and external
service bounds; Part I's authority and proposer-exclusion assumptions.
For the queue/lifecycle pattern, see Ethereum's
\href{https://ethereum.github.io/consensus-specs/specs/fulu/beacon-chain/}{Fulu pending deposits}
and \href{https://ethereum.github.io/consensus-specs/specs/electra/beacon-chain/}{Electra exit and slashing rules}.
Their withdrawal constants are not imported into this protocol.
\end{document}
